%!TEX root = ../../template.tex

\typeout{NT FILE chapter1}

\chapter{Introduction}
\label{cha:introduction}

In this chapter we will contextualize and define the research problem. In Section \ref{ch1:sec:context} we introduce the research context: explaining the concept of the \emph{Semantic Web}, analyzing the industry landscape, enumerate security concerns, that the usage of cloud technologies poses, and show mitigation approaches.  In Section \ref{ch1:sec:problem} we define the problem and research questions. Finally, In Section \ref{ch1:sec:contributions} we present the main contributions from our work.

\section{Context}
\label{ch1:sec:context}

\paragraph{The Semantic Web} In 2001, the term \emph{Semantic Web} \autocite{Berners-Lee2001-ux} was first used to define what would be the evolution of the \gls{WWW}. The traditional \emph{Web} solved the challenge of distributing information at scale and globally, however it was designed for human use. Documents are easily accessible and readable by humans, but for automated agents little more than document retrieval is possible. Contrary to humans, machines cannot reason over the information they access, if they do not have access to the \emph{semantics of the data}.

The \emph{Semantic Web} addresses this issue, by redesigning the representation of information, in a way that not only preserves meaning and relationships but is also machine-readable.
A concept is mapped to a unique identifier, and its meaning is encoded in sets of triples, each triple consisting of a subject, predicate and object (similar to an elementary sentence). \emph{Ontologies} provide a shared understanding of a domain, by formalizing the specification of relevant concepts, and their relations, in such domain.

This novel method of knowledge representation differs from the previous, centralized way, and greatly simplifies the process of interconnecting different knowledge bases. Centralized representations tend to be more expensive to develop, usually needing to be built from scratch \autocite{Neches1991-lx}. Furthermore, for sharing, both parties need to agree on the semantics. In contrast, using ontologies, information is semantically linked by default. In the \emph{Semantic Web}, machines can truly communicate, share and reason about distinct knowledge bases.

Semantic data is specified using a stack of different language specifications. The basic data model is specified using \gls{RDF} \autocite{W3c_undated-vq}, these resources can be organized in hierarchies with \gls{RDFS} \autocite{W3c_undated-mp} and to define ontologies we use \gls{OWL2} \autocite{W3c_undated-dr}. To query semantic data, the standard language is \gls{SPARQL} \autocite{W3c2013-hd, W3c_undated-tw, W3c2008-vf}.


\paragraph{Industry Landscape} Knowledge representation is essential in the industry \autocite{Noy2019-lj} and ontologies help remove ambiguity in terminology.
Nowadays, they are being developed and applied in a variety of domains \autocite{W3c_undated-pd} ranging from the health sector\footnote{The NCI Thesaurus, \autocite{Sioutos2007-ep}, developed by the \gls{NCI}, the \gls{GO} \autocite{Harris2004-tf}, the world’s largest source of information on the functions of genes, or the \gls{SNOMED CT} \autocite{SNOMED-gi}, the most comprehensive, multilingual clinical healthcare terminology in the world, are some existing applications.}, to linguistics \footnote{The WordNet \autocite{Miller1995-ic} is a lexical database for English where related words and concepts are semantically linked.} and even being used by companies like NASA\footnote{Observing the visual representation provided by the Linked Open Data Cloud \autocite{McCrae_undated-or} is a good exercise to perceive the dimension of adoption of this new paradigm.} \autocite{Ashish2005-qz}.

To develop and maintain ontologies, advanced semantic editors are the standard tool for professionals, with Protégé \autocite{Musen2015-kh} being a very popular example. These provide numerous functionalities, usually based on plugins (e.g. automatic verification of errors and consistency in the design, debugging tools, query answering and graphic visualization of data and knowledge).

With datasets gradually growing in dimension, the industry is currently evolving towards the use of collaborative solutions \autocite{Tudorache2008-tu,Tudorache2009-ja} and outsourcing heavy computations, relying on cloud infrastructures or other shared computing platforms. These approaches, albeit inevitable, present new security concerns and engineering challenges.

By 2021, 50 percent of all enterprise workloads were already hosted in \emph{public clouds} \autocite{Security_Compass2021-sh}. As of 2023, worldwide end-user spending on public cloud services is forecast to grow 20.7\% more than the growth forecast for 2022 of 18.8\% \autocite{noauthor_undated-if}.

Research on secure cloud computing \autocite{Carroll2011-me} studies such adoption but also discusses security concerns. Adoption can be explained with the cost efficiency\footnote{Cloud services reduce implementation, maintenance, power, cooling, floor space, storage and operational costs.}, high scalability, flexibility and availability that cloud provides. Additionally, the pay-per-use model is an easy entry point to the ecosystem.

\paragraph{Cloud Security Concerns} While delegation of responsibilities is beneficial for costs, it is prejudicial to control. Using cloud services implies trusting its owner. A \gls{SLA} defines, in accordance to metrics, the service levels of the system that the vendor should provide (e.g. percentage of availability), while also defining penalties in case of not achieving what was agreed-on. However, the risk of loss of privacy and data breaches still exists. A problem that is especially concerning when sensitive data is being stored by such services. This is much more relevant in \emph{public clouds} where, compared to \emph{private clouds}, trust is distributed between various actors and not on a single organization, namely the cloud operator \autocite{Carroll2011-me}.

In cloud storage services, administrators, or hacker with root privileges, might have access to the data \autocite{Bosch2014-eb}. Additionally, these cloud services might outsource computations to third parties whom we do not trust. In this setting, we usually consider the \emph{honest-but-curious}, or \emph{passive}, adversary model. This adversary will follow the protocol properly but keep a record of all its intermediate computations \autocite{Goldreich2001-sr}. Furthermore, it cannot do any other action except attempting to learn private information while observing views of the protocol execution \autocite{Evans2018-it}.

\paragraph{Security Approaches} A first approach to secure data in the cloud would be to use \emph{Encryption at Rest}, which naively encrypts all data stored in the untrusted servers. If a user wants to update or query the data, they will have to download all the data, and decrypt it. This incurs in high network requirements and is computationally expensive on the client-side, hindering all the benefits of using a cloud service to store the data. 

Although we can observe some major database services, like MongoDB, starting to adopt \glsxtrlong{SE} \autocite{noauthor_undated-ds}, the majority of cloud and storage providers only support the former solution.

Currently, active research is being done in the fields of \gls{SE} and \gls{HE}.

\glsxtrlong{HE} is based on the concept of algebra homomorphisms, where both decryption and encryption can be seen as a homomorphism of the relations between the plaintext and the ciphertext \autocite{Acar2018-mp, Fontaine2007-sv}. \glsxtrlong{HE}, based on the type and number of operations supported (e.g. addition or multiplication of encrypted values), can be classified as: \gls{FHE} supporting any number of operations with no restriction on the number of times; \gls{SHE} permitting some types of operations, a limited number of times; \gls{PHE} allowing one type of operation with no restriction on the number of times; and \gls{LFHE} \autocite{Brakerski2014-nm} supporting any type of operations, but with a pre-determined expected depth-$L$ (can be executed up to $L$ times).

\glsxtrlong{SE} uses classical encryption schemes to perform encrypted search queries, sequentially over the ciphertext \autocite{Song2000-hu}, or through the generation of searchable \emph{encrypted indexes} accessible with \emph{tokens}, generated via \emph{trapdoor} functions. \glsxtrlong{SE} can be divided in \gls{SSE}, when it is based on symmetric key cryptography, and \gls{PEKS}, when asymmetric key algorithms are used. \gls{SSE} is often more practical than \gls{PEKS}, as symmetric cryptographic schemes tend to be more efficient than their asymmetric counterparts \autocite{Pearson2017-ve}. These techniques may leak some information about the \emph{indexes}, \emph{search pattern}, and \emph{access pattern} during execution. Most approaches leak, at least, the \emph{search} and \emph{access patterns} \autocite{Bosch2014-eb}. Recently, to mitigate information leakage, \emph{hybrid} solutions combining the software approach of \gls{SE} with \emph{trusted hardware} have been proposed \autocite{Ferreira2020-ze}.

Besides \gls{HE} and \gls{SE}, \gls{ORAM} \autocite{Goldreich1996-yb} could theoretically be used as an alternative. However, \gls{ORAM} and \gls{HE} continue to be too inefficient for real world applications, when considered as an exclusive solution \autocite{Bosch2014-eb}.

\section{Problem}
\label{ch1:sec:problem}

Usually, ontologies only contain public, and non-sensitive data, but sometimes they are used alongside electronic health records and clinical decision systems. These systems heavily manipulate personal data, which is sensitive. Additionally, ontologies themselves have an inherent private value, due to the investment of time and money in their development. With all these potential security needs, a solution that supports the secure development, sharing and usage of ontologies alongside sensitive data is essential.

Taking into consideration the efficiency, security trade-offs and ease of implementation from the various alternative techniques, a pure software \gls{SSE} approach seems to be the most practical and easier to adopt solution, not requiring excessive computational power or special trusted hardware. Recently, Chase and Kamara, on their work on \gls{STE}, generalized index-based \gls{SSE} to support queries on \emph{structured} data \autocite{Chase2010-xx}.

Semantic data can be represented as a graph and \gls{STE} supports various types of query operations in graph-structured data. Some existing schemes, adapted to conceptual graphs, need to pre-compute all possible query answers \emph{a priori}, which is not practical for large datasets \autocite{Poh2012-qs}. Furthermore, there can be leakage of information about the topology of the data while trying to pre-compute all query answers. 

Recent developments do not require pre-computation of all query answers. The \emph{CryptGraphDB} \cite{Aburawi2018-pq} system is an adaptation of CryptDB \cite{Popa2011-zb, Popa2012-bw, Popa2011-ci} capable of answering queries over conceptual graphs, but it was developed for the cypher \cite{Neo4j_Inc_undated-ft} language and not \gls{SPARQL}.

Another work, focused on access control for RDF datasets, did also implement a solution for retrieving simple triple patterns over an encrypted dataset, but, due to performance needs, hashes were used for indexing the encrypted triples. This leaks the structure of the dataset, which they tried to mitigate by creating dummy triples \cite{Fernandez2017-bh}.

The \gls{GOOSE} \cite{Ciucanu2020-lx} framework, demonstrates how to use multiparty computation techniques to answer \gls{UCRPQ} over encrypted graphs, only relying on a standard SPARQL engine. But deterministic data obfuscation techniques are used, and therefore the structure of the graph is also leaked. This framework only focuses on SPARQL read queries.

Although, a lot of progress has been made in this field of research there are still various open questions. In this thesis, we will focus on answering the following questions:

\begin{enumerate}
    \item \emph{How to support secure and expressive read and write \gls{SPARQL} queries over encrypted data?}
    
    A practical solution ought to support more than basic triple pattern retrieval, it must be capable of querying complex patterns. As such it should be able to execute secure set operations without the client needing to decrypt the intermediate values (obtained from simple triple patterns). Furthermore, it must be able to execute write queries, so that datasets can be updated. 

    \item \emph{How can we add reasoning capabilities to the solution?}
    
    A simple \gls{SPARQL} engine does not possess reasoning capabilities. The query results may or may not possess inferred values, depending on the underlying storage solution. When working with ontologies, a system should have reasoning capabilities to take full advantage of the semantics preserved within the data.

    \item \emph{How can we minimize the information leakage and enforce access control?}
    
    A data outsourcing solution should support granular access control over the shared datasets. \gls{SSE} solutions always leak some information about the data, we will investigate how to minimize such leakage during query execution.

\end{enumerate}


\section{Contributions}
\label{ch1:sec:contributions}

In this thesis we aim to develop a practical and secure solution for the collaborative development, distribution and usage of semantic data knowledge bases. Our main contributions are the following:


\begin{description}
    \item \textbf{SPARQL-EVAL \& SPARQL-EVAL-II - Distributed two-party \gls{SPARQL} query evalua-}\linebreak \textbf{tion protocol:} We developed two \gls{SPARQL} query evaluation protocols which support a subset of queries: queries target solely the default graph; the allowed modifier are DISTINCT, ORDERBY and LIMIT; graph patterns must not contain property paths, nested subqueries queries, graph selection, nor triple patterns composed solely by variables. In SPARQL-EVAL-II we use query rewriting techniques to retain reasoning capabilities. 
    
    \item \textbf{SSE-TP - A distributed two-party triple pattern \gls{SSE} protocol:} 
    SSE-TP is protocol for encrypting \gls{RDF} data in a way that the data remains searchable via the specification of triple patterns. 

    \item \textbf{SSE-SPARQL-EVAL - A distributed three-party \gls{SPARQL} query evaluation protocol} \linebreak \textbf{ for SSE-TP encrypted triplestores:}   
    SSE-SPARQL-EVAL is an adaptation of the SPARQL-EVAL-II protocol that works with SSE-TP encrypted triplestores. 

    \item \textbf{\gls{SEnTriS} - A multiparty searchable encrypted triplestore sharing protocol:}
    
    \gls{SEnTriS} is an adaptation of the SSE-SPARQL-EVAL protocol for a multiuser setting. Encrypted triplestores are shareable, remain queryable, possess reasoning capabilities and support role based access control. Access to triplestores is controlled with the usage of ephemeral access tokens. Concurrent writes are resolved with a lock-based strategy. Secrets are distributed among user with access to the triplestores. The protocol is implemented in two versions:
    \begin{itemize} 
    
        \item \emph{\gls{SEnTriS} V1}: Data encryption retains elements from \emph{CryptDB} \autocite{Popa2011-zb, Popa2012-bw, Popa2011-ci}. Subjects, predicates and objects are encrypted in two layers. The outer layer is probabilistically encrypted, whereas the second is deterministically encrypted. When equality checks are required, we peel the outer layer of the encryption onions and compare the values.
    
        \item \emph{\gls{SEnTriS} V2}: Data is encrypted using the \gls{DGK} cryptosystem \autocite{Damgard2007-km, Damgard2009-my, Franz2011-sm, Makkes2010-qb}, an additive homomorphic scheme developed for secure integer comparison. The value of a subject, predicate or object is mapped to an \emph{equality tags}. Trapdoors index the \emph{equality tags} and equality tags reference the probabilistically encrypted subjects, predicates and objects. During query evaluation, we mask the values of the equality tags, via re-encryption. When doing equality checks, we decrypt the equality tags and compare the values. These remain deterministic, but different for every query evaluation.
    \end{itemize}
\end{description}

\section{Document Organization}
\label{ch1:sec:organization}

The remaining of the document will be organized as follows:

\begin{description}
    \item [Chapter 2] discusses related work on \emph{Semantic Web} and \emph{Secure Computations} and compares existing solutions for secure semantic data outsourcing and \gls{SPARQL} evaluation.
    \item [Chapter 3] presents the solution.
    \item [Chapter 4] describes the implementation details of each of the system's components.
    \item [Chapter 5] presents the experimental evaluation procedures, describes and discusses the obtained results.
    \item [Chapter 6] summarizes main conclusions from the research and suggests future work.
\end{description}